\documentclass[14pt,a4paper,oneside]{extbook}
\usepackage[utf8]{vietnam}
\usepackage{amsmath,amssymb,mathrsfs,fancyhdr,enumerate,multirow,makecell}
\usepackage{tkz-euclide,tikz-3dplot,tikz,tkz-tab}
\usetikzlibrary{arrows.meta,calc,intersections,angles,quotes,backgrounds,shapes.geometric,patterns,through,shadows,positioning}
\usepackage{yhmath}
\usepackage[top=0.7cm, bottom=0.7cm, left=1.0cm, right=1.0cm, headheight=16pt, headsep=3mm, footskip=5mm]{geometry}
\usepackage[hidelinks,unicode]{hyperref}
\usepackage{qrcode}
\usepackage[many]{tcolorbox}
\usepackage{fontawesome}

\pagestyle{fancy}
\fancyhf{}
\fancyhead[L]{\footnotesize\bfseries\color{teal!85!black}\faBook~SỔ TAY KIẾN THỨC TOÁN 5}
\fancyhead[R]{\footnotesize\bfseries\color{blue!85!black}Lớp Toán Thầy Lĩnh -- loptoanthaylinh.vn}
\fancyfoot[C]{\footnotesize\bfseries\color{gray!80!black}-- Sổ tay kiến thức: Bài 33. Ôn tập diện tích, chu vi một số hình phẳng --}
\renewcommand{\headrulewidth}{0.4pt}

\begin{document}

% --- BANNER TIÊU ĐỀ BÀI HỌC ---
\noindent
\begin{tcolorbox}[
  enhanced,
  colback=teal!8!white,
  colframe=teal!80!black,
  arc=3mm,
  boxrule=1.1pt,
  width=\linewidth,
  halign=center,
  top=2.2mm, bottom=2.2mm,
  left=3mm, right=3mm,
  drop fuzzy shadow
]
  {\large\bfseries\color{teal!95!black}\faBookmark\ BÀI 33. ÔN TẬP DIỆN TÍCH, CHU VI MỘT SỐ HÌNH PHẲNG}\\[2pt]
  {\scriptsize\color{black!75}\itshape Bảng tra cứu toàn bộ công thức Chu vi và Diện tích từ lớp 1 đến lớp 5 trên 1 trang A4}
\end{tcolorbox}

\vspace{3.2mm}
% --- KHỐI 1: BẢNG TRA CỨU TOÀN TẬP CHU VI VÀ DIỆN TÍCH ---
\noindent
\begin{tcolorbox}[
  enhanced,
  title={\footnotesize\bfseries\color{white}\faTable\ 1. Bảng tra cứu toàn diện công thức Chu vi và Diện tích},
  colback=blue!2!white,
  colframe=blue!75!black,
  arc=2.5mm, boxrule=0.8pt, top=1.8mm, bottom=1.8mm, left=2mm, right=2mm
]
  \scriptsize
  \centering
  \begin{tabular}{|l|c|c|l|}
    \hline
    \textbf{Hình học} & \textbf{Chu vi ($P$ hoặc $C$)} & \textbf{Diện tích ($S$)} & \textbf{Công thức biến đổi quan trọng} \\ \hline
    \textbf{Hình vuông} & $P = a \times 4$ & $S = a \times a$ & Cạnh $a = P : 4$ \\ \hline
    \textbf{Hình chữ nhật} & $P = (a + b) \times 2$ & $S = a \times b$ & Nửa chu vi $= a + b = P : 2$; $a = S : b$ \\ \hline
    \textbf{Hình bình hành} & $P = (a + b) \times 2$ & $S = a \times h$ & Chiều cao $h = S : a$; Đáy $a = S : h$ \\ \hline
    \textbf{Hình thoi} & $P = a \times 4$ & $S = \dfrac{m \times n}{2}$ & Tích 2 đường chéo: $m \times n = S \times 2$ \\ \hline
    \textbf{Hình tam giác} & $P = a + b + c$ & $S = \dfrac{a \times h}{2}$ & $h = \dfrac{2S}{a}$; Đáy $a = \dfrac{2S}{h}$ \\ \hline
    \textbf{Hình thang} & $P = a + b + c + d$ & $S = \dfrac{(a + b) \times h}{2}$ & $h = \dfrac{2S}{a + b}$; Tổng 2 đáy $= a + b = \dfrac{2S}{h}$ \\ \hline
    \textbf{Hình tròn} & $C = d \times 3,14 = 2r \times 3,14$ & $S = r \times r \times 3,14$ & Đường kính $d = C : 3,14$; $r \times r = S : 3,14$ \\ \hline
  \end{tabular}
\end{tcolorbox}

\vspace{3mm}

% --- KHỐI 2: DẠNG TOÁN TỈ LỆ VÀ BIẾN ĐỔI KÍCH THƯỚC ---
\noindent
\begin{minipage}[t]{0.52\textwidth}
  \begin{tcolorbox}[
    enhanced,
    title={\footnotesize\bfseries\color{white}\faSliders\ 2. Quy luật tỉ lệ kích thước và diện tích},
    colback=teal!2!white,
    colframe=teal!75!black,
    arc=2.5mm, boxrule=0.8pt, top=1.8mm, bottom=1.8mm, left=2.2mm, right=2.2mm
  ]
    \footnotesize
    $\bullet$ Nếu một kích thước tăng $k$ lần (kích thước kia giữ nguyên): Diện tích tăng \textbf{$k$ lần}.\\
    $\bullet$ Nếu cả 2 kích thước đều tăng $k$ lần: Diện tích tăng \textbf{$k \times k$ lần}.\\
    \textit{VD:} Bán kính hình tròn tăng gấp $3$ lần $\rightarrow$ Chu vi tăng $3$ lần, nhưng Diện tích tăng $3 \times 3 = \mathbf{9\text{ lần}}$.
  \end{tcolorbox}
\end{minipage}%
\hfill
\begin{minipage}[t]{0.46\textwidth}
  \begin{tcolorbox}[
    enhanced,
    title={\footnotesize\bfseries\color{white}\faCheckSquareO\ 3. Phương pháp giải hình phức hợp},
    colback=purple!2!white,
    colframe=purple!70!black,
    arc=2.5mm, boxrule=0.8pt, top=1.8mm, bottom=1.8mm, left=2mm, right=2mm
  ]
    \footnotesize
    $\bullet$ \textbf{Cộng dồn}: Tách hình thành các mảnh chữ nhật, tam giác, nửa hình tròn $\rightarrow$ tính rời rồi cộng lại.\\
    $\bullet$ \textbf{Trừ phần khuyết}: Lấy diện tích khung bao trừ đi phần không thuộc hình cần tính.
  \end{tcolorbox}
\end{minipage}

\vspace{3mm}

% --- KHỐI 3: SƠ ĐỒ TƯ DUY ---
\noindent
\begin{tcolorbox}[
  enhanced,
  title={\footnotesize\bfseries\color{white}\faSitemap\ Sơ đồ tư duy: Chiến lược giải toán chu vi & diện tích},
  colback=yellow!2!white,
  colframe=teal!85!black,
  arc=3mm, boxrule=0.9pt, top=1.8mm, bottom=1.8mm, left=1.5mm, right=1.5mm,
  halign=center
]
  \begin{tikzpicture}[
    scale=0.92,
    every node/.style={transform shape},
    b1/.style={rectangle, rounded corners=4pt, fill=blue!15, draw=blue!70!black, text=blue!90!black, font=\bfseries\scriptsize, inner sep=4pt, align=center},
    b2/.style={rectangle, rounded corners=4pt, fill=teal!15, draw=teal!70!black, text=teal!90!black, font=\bfseries\scriptsize, inner sep=4pt, align=center},
    b3/.style={rectangle, rounded corners=4pt, fill=green!15, draw=green!70!black, text=green!90!black, font=\bfseries\scriptsize, inner sep=4pt, align=center}
  ]
    \node[b1] (n1) at (-5.8, 0) {1. Đồng nhất đơn vị đo\\(cm, dm, m, ha, km)};
    \node[b2] (n2) at (0, 0) {2. Nhận diện hình\\Chọn đúng công thức};
    \node[b3] (n3) at (5.8, 0) {3. Tính toán \& Ghi đơn vị\\(Độ dài vs Diện tích)};

    \draw[teal!85!black, thick, -{Stealth[length=2.5mm]}] (n1) -- (n2);
    \draw[teal!85!black, thick, -{Stealth[length=2.5mm]}] (n2) -- (n3);
  \end{tikzpicture}
\end{tcolorbox}

\vspace{3mm}

% --- KHỐI 4: KỸ NĂNG TÍNH NHANH & BẪY TRẮC NGHIỆM ---
\noindent
\begin{minipage}[t]{0.54\textwidth}
  \begin{tcolorbox}[
    enhanced,
    title={\footnotesize\bfseries\color{white}\faBolt\ Mẹo nhớ mẫu số chia 2},
    colback=green!3!white,
    colframe=green!60!black,
    arc=2.5mm, boxrule=0.8pt, top=1.8mm, bottom=1.8mm, left=2.2mm, right=2.2mm
  ]
    \footnotesize
    $\bullet$ Các hình có chia 2 trong công thức diện tích:\\
    $\triangleright$ \textbf{Hình tam giác}: $(a \times h) : 2$.\\
    $\triangleright$ \textbf{Hình thoi}: $(m \times n) : 2$.\\
    $\triangleright$ \textbf{Hình thang}: $(a + b) \times h : 2$.\\
    Ghi nhớ: Cứ hình nào có ``đáy nhọn'' hoặc ``vát nghiêng'' là phải chia đôi!
  \end{tcolorbox}
\end{minipage}%
\hfill
\begin{minipage}[t]{0.44\textwidth}
  \begin{tcolorbox}[
    enhanced,
    title={\footnotesize\bfseries\color{white}\faExclamationTriangle\ Mẹo tránh bẫy trắc nghiệm},
    colback=red!2!white,
    colframe=red!70!black,
    arc=2.5mm, boxrule=0.8pt, top=1.8mm, bottom=1.8mm, left=2.2mm, right=2.2mm
  ]
    \footnotesize
    $\bullet$ Đơn vị chu vi là đơn vị độ dài ($\text{m}, \text{cm}$); đơn vị diện tích là đơn vị vuông ($\text{m}^2, \text{cm}^2$). Viết nhầm mũ $2$ là bị gạch điểm!\\
    $\bullet$ Đọc kĩ đề: Cho chu vi mà hỏi cạnh hay hỏi diện tích.
  \end{tcolorbox}
\end{minipage}

\end{document}
